%% kapa: exam=1 round=25 year=2014 period=1 form=A subject=civil_law subject_ko=민법 seq=25 number=25
%% source: Q-Net 감정평가사 기출문제 (cst00309) artlSeq=5101207 file=extracted/1008/5101207/2014년 제25회 감정평가사 1차시험 기출문제/1교시(A형).hwp
%% conversion: kordoc 4.24.51 -> tools/gen1.py (auto); review=proofread-sonnet; flags=none
\begin{question}{25}
소멸시효에 관한 설명으로 옳지 않은 것은? (다툼이 있으면 판례에 의함)
\begin{choices}
  \item 소멸시효는 그 기산일에 소급하여 효력이 생긴다.
  \item 소멸시효이익은 시효가 완성되기 전에 미리 포기하지 못한다.
  \item 주된 권리가 시효로 소멸하면 종된 권리에도 그 효력이 미친다.
  \item 20년의 시효에 걸리는 권리에 확정판결이 있으면, 그 시효기간은 10년으로 단축된다.
  \item 시효는 법률행위에 의하여 이를 가중할 수 없으나 경감할 수는 있다.
\end{choices}
\end{question}
\answer{4}
